\subsection{Tests und Builds}
\label{subsec:tests-builds}

Die acht Testsuiten decken Tooling, Schema, API, Datenbank, PostgreSQL,
Frontend, E2E und Tauri/Rust ab. \texttt{test --suite all} bildet den
vollständigen projektbezogenen Einstieg; Berichte sind optional
\parencite{kleiveist2026control,kleiveist2026tooling}.
Die Tooling-Suite prüft unter anderem Generator, Profile und Konfiguration; die
Komponentensuiten delegieren an Pytest, Vitest, Playwright oder Cargo. Damit
bleiben Testbereiche getrennt ausführbar.

Die Auswahl bleibt profilabhängig: Fehlt ein Feature, meldet die Suite
\texttt{SKIP}; fehlen für ein aktives Feature Quellen oder Voraussetzungen,
meldet sie \texttt{FAIL}. PostgreSQL ohne Test-URL und E2E ohne
Playwright-Konfiguration werden übersprungen. Der Runner kann fehlende
Backend-Umgebungen reparieren und für E2E die Entwicklungsdienste verwalten.
Diese Semantik unterscheidet bewusst fehlende Features von defekten aktiven
Komponenten. Ein übersprungener Test ist deshalb kein Erfolgsnachweis.

Die Buildpfade bleiben getrennt: Web erzeugt Vite-Ausgabe und ZIP, Desktop
delegiert Ziel und Bundle an Tauri, Container baut für Cloudprofile Frontend-,
Backend- oder beide Images. \texttt{container validate} prüft das
Compose-Modell ohne Start \parencite{kleiveist2026container,kleiveist2026release}.
Der Webpfad legt zusätzlich ein ZIP-Paket ab. Desktop-Builds bleiben an ein
aktives Tauri-Feature und die jeweilige Zielplattform gebunden.

GitHub Actions verwendet dieselben öffentlichen Einstiege für Core,
Profilmatrizen, PostgreSQL und Desktopziele. CI ergänzt dafür Betriebssysteme,
Laufzeiten und einen temporären PostgreSQL-Dienst
\parencite{kleiveist2026ci}.
Lokale und externe Abläufe nutzen damit vergleichbare Einstiegspunkte, laufen
aber nicht unter identischen Umgebungsbedingungen.
\beginnerinput{workspace/statement/500_tooling_workflow/550}
